\begin{longtable}{@{}P{.07\linewidth}P{.64\linewidth}P{.22\linewidth}@{}}
\caption{Afgeleide requirements. MOET betekent vereist binnen het voorgestelde ontwerp, niet zelfstandig wettelijk verplicht.}\label{tab:req}\\
\toprule Eis & Voorgestelde ontwerpverplichting & Grondslag \\ \midrule\endfirsthead
\toprule Eis & Voorgestelde ontwerpverplichting & Grondslag \\ \midrule\endhead
\hypertarget{R01}{}R01 & Typeer term, begrip, definitie, klasse, modelelement, record en product afzonderlijk. & \citep{S11,S13,S31,S08} \\
\hypertarget{R02}{}R02 & Verbind een definitie met bron, context, versie, verantwoordelijke en toepassingsperiode. & \citep{S14,S25,S33} \\
\hypertarget{R03}{}R03 & Bewaar identifier-scope en onderscheid entiteit-ID, record-ID, versie-ID en document-URL. & \citep{S28,S38,S39,L08} \\
\hypertarget{R04}{}R04 & Koppel data-elementen aan begrip, eigenschap en waarde-/conceptueel domein. & \citep{S08,S13} \\
\hypertarget{R05}{}R05 & Maak syntax, semantic encoding, dataspecificatie en decoderingsinformatie controleerbaar. & \citep{S03} \\
\hypertarget{R06}{}R06 & Onderscheid geldigheid, registratie, publicatie en eventtijd; bewaar correctiehistorie. & \citep{L14,S37,S22} \\
\hypertarget{R07}{}R07 & Leg bronrecordversie, afleidingsactiviteit, verantwoordelijke en gebruikte modelversies vast. & \citep{S04,S33} \\
\hypertarget{R08}{}R08 & Scheid wettelijke authenticiteit, beheerautoriteit, feitelijke juistheid en technische afzenderidentiteit. & \citep{S25,S21,S33} \\
\hypertarget{R09}{}R09 & Publiceer kwaliteit met metriek, populatie, periode, meetmethode, streefwaarde en resultaat. & \citep{S05,S06,S12,S34,S48,L15} \\
\hypertarget{R10}{}R10 & Versieer definities, ontologie, model, shape, context, mapping, dataset en interface onafhankelijk maar samenhangend. & \citep{S14,S23,S47,S18} \\
\hypertarget{R11}{}R11 & Maak mappings gericht, getypeerd, context- en versiegebonden; motiveer equivalentie. & \citep{S31,S30,S13,S08} \\
\hypertarget{R12}{}R12 & Registreer voorstel, review, vaststelling, publicatie, deprecatie en intrekking met verantwoordelijke. & \citep{S10,S23,L07,L17} \\
\hypertarget{R13}{}R13 & Beschrijf product-ID, doelgroep, doel, eigenaar, datasets, diensten, kwaliteit, gebruikscondities en afhankelijkheden. & \citep{L12,S16,S47,S27} \\
\hypertarget{R14}{}R14 & Valideer uitwisselprofielen en behoud betekenis bij API- en RDF/JSON-LD-representaties. & \citep{S18,S19,S20,S35} \\
\hypertarget{R15}{}R15 & Verbind notificatie met bronobject en schema; leg verwerking van duplicaten, volgorde en herstel vast. & \citep{S22,L14} \\
\hypertarget{R16}{}R16 & Scheid logische inferentie, constraints op beschikbare data en externe werkelijkheidscontrole. & \citep{S29,S30,S32,S12} \\
\hypertarget{R17}{}R17 & Bied machinevindbare betekenis en provenance en markeer ontbrekend of betwist bewijs. & \citep{L11,S33,L13,L16} \\
\hypertarget{R18}{}R18 & Beoordeel FDS-toelating met afzonderlijke verplichte en aanbevolen afspraken. & \citep{S27} \\
\hypertarget{R19}{}R19 & Ondersteun terugmelding met bronhouder, status, afhandeling en traceerbare correctie. & \citep{S25,S26,S27} \\
\hypertarget{R20}{}R20 & Koppel bewaarde informatieobjecten en versies aan een afzonderlijk bewaarmetadata-profiel. & \citep{S24,S15} \\
\hypertarget{R21}{}R21 & Maak toepasselijkheid van juridische en standaardverplichtingen afzonderlijk toetsbaar. & \citep{S42,S43,S25} \\
\hypertarget{R22}{}R22 & Toets identity, role en phase met ontologische analyse en vergelijk de benodigde inspanning. & \citep{L08,L09,L10,S40,S41} \\
\hypertarget{R23}{}R23 & Maak uitspraaksoort, registratiehandeling, productiemethode, onzekerheidsmodel, gezagsgrond en bewijsstatus afzonderlijk machineleesbaar per bewering. & \citep{S25,S33,S04,L11,L16} \\
\bottomrule\end{longtable}